% Chapter-local TikZ diagram; final width about 168 mm, no external assets.
% Solid arrows carry records. The dashed return arrow requests evidence checks.
\begin{tikzpicture}[
  x=1mm,y=1mm,
  font=\figurefont\fontsize{9}{11}\selectfont,
  box/.style={rectangle,align=left,inner sep=2mm,line width=0.85pt,
    draw=BookRule,fill=white,text width=42mm,minimum height=19mm},
  input/.style={box,draw=FigureInput,fill=FigureInput!8},
  model/.style={box,draw=FigureModel,fill=FigureModel!8},
  output/.style={box,draw=FigureOutput,fill=FigureOutput!8},
  flow/.style={line width=1pt,draw=BookInk,-{Latex[length=2mm]}},
  feedback/.style={line width=0.85pt,draw=FigureModel,dashed,
    -{Latex[length=2mm]}},
  heading/.style={font=\figurefont\bfseries\fontsize{9}{11}\selectfont,
    text=BookInk,anchor=west}
]
  \node[heading] at (0,17) {A\quad 每条事实保留可回查的依据};
  \node[input] (text) at (23,0) {
    \textbf{原始文本}\\
    文档ID、版本、发布日期\\
    保留未经覆盖的原文
  };
  \node[input] (mention) at (84,0) {
    \textbf{提及定位}\\
    字符跨度、类型\\
    词元到原文的映射
  };
  \node[model] (binding) at (145,0) {
    \textbf{对象绑定}\\
    共指簇、实体身份\\
    活动与角色归属
  };
  \node[model] (candidate) at (145,-34) {
    \textbf{关系或事件候选}\\
    对象、角色、极性\\
    时间、条件、证据句
  };
  \node[model] (check) at (84,-34) {
    \textbf{记录核验}\\
    模式、身份、原文支持\\
    限定条件与来源
  };
  \node[output] (record) at (23,-34) {
    \textbf{带来源记录}\\
    规范键、字段与值\\
    有效时间、写入时间
  };
  \draw[flow] (text.east) -- (mention.west);
  \draw[flow] (mention.east) -- (binding.west);
  \draw[flow] (binding.south) -- (candidate.north);
  \draw[flow] (candidate.west) -- (check.east);
  \draw[flow] (check.west) -- (record.east);
  \draw[feedback] (check.north) -- node[right,text=FigureModel] {回查} (mention.south);
  \draw[BookRule,line width=0.7pt] (0,-52) -- (168,-52);
  \node[heading] at (0,-60) {B\quad 同一活动改期：保留旧值，不把两次主张当成两场活动};
  \node[input,text width=75mm,minimum height=21mm] (old) at (40,-80) {
    \textbf{旧通知：4月14日发布}\\
    活动编号 ML-0418\\
    原排期为4月17日，未给时刻
  };
  \node[input,text width=75mm,minimum height=21mm] (new) at (127,-80) {
    \textbf{更正通知：4月15日发布}\\
    同一活动改至4月18日14:00\\
    “原定”与“现改至”连接新旧值
  };
  \node[output,text width=163mm,minimum height=25mm] (versions) at (84,-119) {
    \textbf{同一活动键，两个排期版本（年份均为2025年）}\\
    旧值：4月17日；排期主张有效期 $[\text{4月14日},\text{4月15日})$；保留旧来源\\
    新值：4月18日14:00；排期主张自4月15日起有效；关联更正来源\\
    系统若4月16日才读到更正，写入日为4月16日，有效起日仍为4月15日
  };
  \draw[BookInk,line width=1pt] (old.south) -- (40,-98) -- (84,-98);
  \draw[BookInk,line width=1pt] (new.south) -- (127,-98) -- (84,-98);
  \draw[flow] (84,-98) -- (versions.north);
\end{tikzpicture}
